%
% graph.tex -- Template für standalone TIKZ Bilder
%
% (c) 2019 Prof Dr Andreas Müller, Hochschule Rapperswil
%
\documentclass[tikz,12pt]{standalone}
\usepackage{amsmath}
\usepackage{times}
\usepackage{txfonts}
\usepackage{pgfplots}
\usepackage{csvsimple}
\definecolor{darkred}{rgb}{0.8,0,0}
\definecolor{darkgreen}{rgb}{0,0.6,0}
\usetikzlibrary{arrows,intersections,math,calc}
\begin{document}
\def\skala{4}
\begin{tikzpicture}[>=latex,thick,scale=\skala]

\begin{scope}
\clip (-0.5,0) rectangle (2.5,1.5);
\fill[color=blue!20] 
	plot[domain=0:1,samples=30] ({\x},{\x^2})
	--
	plot[domain=0:2,samples=30] ({1+\x},{exp(-\x)})
	--
	(3,0)
	--
	cycle;
\end{scope}
\draw[->] (-0.6,0) -- (2.6,0) coordinate[label={$x$}];
\draw[->] (0,-0.1) -- (0,1.6) coordinate[label={right:$y$}];
\begin{scope}
\clip (-0.5,-0.01) rectangle (2.5,1.5);
\foreach \c in {-2,-1.8,...,2.4}{
	\draw[color=darkred,line width=1.2pt]
		plot[domain=0:1.3,samples=30]
			({\x+\c},{\x^2});
}
\foreach \d in {-4,-3.8,...,3}{
	\draw[color=darkgreen,line width=1.2pt]
		plot[domain=-1:3,samples=30]
			({\x+\d},{exp(-\x)});
}
\draw[color=orange,line width=1.4pt] (0,0) -- (2.5,0);
\end{scope}
\fill[color=blue] (1,1) circle[radius=0.02];
\node[color=blue] at (1,1) [left] {$P$};
\draw (1,-0.02) -- (1,0.02);
\node at (1,0) [below] {$1$};
\draw (2,-0.02) -- (2,0.02);
\node at (2,0) [below] {$2$};
\draw (-0.02,1) -- (0.02,1);
\node at (0,1) [left] {$1$};

\end{tikzpicture}
\end{document}

